\section{Related Work}
\subsection{AI Researcher Frameworks}

Recent works have introduced end-to-end (E2E) frameworks designed to fully automate the scientific research lifecycle. The AI Scientist-v1~\citep{lu2024aiscientist} introduced automated experimentation and drafting via code templates, while v2~\citep{yamada2025aiscientistv2} increases autonomy using agentic tree-search. Concurrent works expand on these paradigms through distinct technical approaches: Cycle Researcher~\citep{weng2024cycleresearcher} employs an iterative refinement loop with reinforcement learning from reviewer feedback, while OmniScientist~\citep{shao2025omniscientist} couples an agentic framework with structured knowledge and tool-based reasoning. Other systems, such as InternAgent~\citep{team2025internagent} and AI Co-Scientist~\citep{gottweis2025towards}, support human-in-the-loop collaboration in domain-specific applications. While effective at autonomous execution, the writing modules in these frameworks are rigidly coupled to their internal experimental loops. Consequently, they cannot directly function as standalone pipelines to synthesize unconstrained, human-provided pre-writing materials into submission-ready manuscripts.

\begin{table}[tbp]
    \centering
    \resizebox{\linewidth}{!}{%
    \setlength{\tabcolsep}{4pt}
    \begin{tabular}{@{} l c c c c c @{}}
        \toprule
        \textbf{Method} & 
        \textbf{\begin{tabular}{@{}c@{}}E2E \mylatex \\ Manuscript Gen.\end{tabular}} & 
        \textbf{\begin{tabular}{@{}c@{}}Decoupled \\ Standalone Writer\end{tabular}} & 
        \textbf{\begin{tabular}{@{}c@{}}Robust to \\ Unstructured Input\end{tabular}} & 
        \textbf{\begin{tabular}{@{}c@{}}Targeted Lit \\ Review Gen.\end{tabular}} & 
        \textbf{\begin{tabular}{@{}c@{}}Conceptual \\ Diagram Gen.\end{tabular}} \\
        \midrule
        PaperRobot \citep{wang2019paperrobot}            & \xmark & \xmark & \xmark & \xmark & \xmark \\
         data-to-paper \citep{technion2024datatopaper}    & \xmark & \cmark & \cmark & \xmark & \xmark \\
        % XtraGPT \citep{chen2025xtragpt}                  & \xmark & \cmark & \cmark & \xmark & \xmark \\
        AI-Researcher \citep{tang2025ai}                 & \cmark & \xmark & \xmark & \cmark & \xmark \\
        Cycle Researcher \citep{weng2024cycleresearcher} & \cmark & \xmark & \xmark & \xmark & \xmark \\
        AI Scientist-v2 \citep{yamada2025aiscientistv2}  & \cmark & \xmark & \cmark & \cmark & \xmark \\
        \ourmethod\ \textbf{(Ours)}                      & \cmark & \cmark & \cmark & \cmark & \cmark \\
        \bottomrule
    \end{tabular}
    }
    \caption{Comparison of \ourmethod{} against existing AI writing systems.}
    \label{tab:comparison}
\end{table}

\subsection{Automated Writing and Literature Synthesis}
Early automated writing systems, such as PaperRobot~\citep{wang2019paperrobot}, generated incremental text sequences conditioned on entity relation graphs but lacked the capacity to synthesize complex, data-driven scientific narratives. More recent LLM-based writing assistants, including Prism~\citep{openai_prism_2026}, excel at stylistic refinement and local text editing, but typically rely on structured inputs or human guidance and do not natively support end-to-end manuscript generation from raw experimental logs. Similarly, data-to-paper~\citep{technion2024datatopaper} focuses on translating structured analytical results into text with backward traceability. However, it primarily operates on structured, well-defined inputs and is less suited to flexibly drafting from early-stage, under-specified raw materials.

Beyond constructing a data-driven narrative from experimental results, synthesizing prior literature remains a distinct and critical challenge in automated authoring. Recent work targets this problem via automated literature review generation. For instance, AutoSurvey2~\citep{wu2025autosurvey2} uses a multi-stage retrieval-augmented pipeline to generate comprehensive surveys, while SurveyGen-I~\citep{chen2025surveygen} focuses on iteratively refining topic outlines. Similarly, LiRA~\citep{go2025lira} employs a multi-agent workflow that structures, drafts, and refines literature reviews from retrieved papers. While highly effective at generating long-form surveys, these systems are not designed for writing targeted related work sections. Consequently, they often lack the contextual awareness needed to selectively contrast prior work and clearly motivate the specific research gap of a new method. 

As summarized in Table~\ref{tab:comparison}, existing pipelines show key limitations when used as standalone manuscript generators. Frameworks like CycleResearcher~\citep{weng2024cycleresearcher} require a structured BibTeX reference list as input, which is rarely available at the start of writing, and fail on unstructured inputs outside this format. While AI-Researcher~\citep{tang2025ai} and AI Scientist-v2~\citep{yamada2025aiscientistv2} accept unstructured topic descriptions, their writing modules rely on artifacts produced earlier in their internal pipelines and do not generate conceptual scientific diagrams. In contrast, \ourmethod{} can process unconstrained pre-writing materials without relying on pre-existing plots or ground-truth references. By orchestrating specialized agents to autonomously synthesize targeted literature, generate comprehensive visuals, and craft coherent scientific narratives, \ourmethod{} establishes a robust, end-to-end pipeline for AI research manuscript generation.